\cotefile{157-3}
\scanpage{50}
XII

\unsure{Caractérisation} de $W_\infty$ etc

Sommes \unsure{amalgamées} et carrés $W$-cocartésiens.

Théorie des carrés \margin{\unsure{$\downarrow$|l}} \unsure{carrés} \inserted{$W$-}cartésiens et \deleted{\unsure{$W$}}$W$-cocartésiens dans les catégories de \unsure{modèles}

\figure{\begin{tikzpicture}[x=1cm,y=1cm]
\draw[line width=.2pt] (-3.56,1.49) -- (-2.86,1.53);
\draw[line width=.2pt] (-3.08,1.12) .. controls (-2.83,.91) and (-2.52,.65) .. (-2.48,.93) -- (-2.50,1.04);
\draw[line width=.35pt] (0.47,.35) .. controls (.30,.31) and (.08,.32) .. (.02,.40) .. controls (-.04,.47) and (.10,.57) .. (.13,.54) -- (.35,.37) .. controls (.62,.22) and (1.05,.06) .. (1.46,0) .. controls (1.51,-.02) and (1.50,.05) .. (1.47,.09);
\end{tikzpicture}}

% archive sheet number: 50
\scanpage{51}
Caractérisation de $W_\infty$ (\unsure{Hot}-équivalences) dans $\mathrm{Cat}$

$\overline{\underline{\mathrm{XII}}}$

\unsure{Axiomes} $W(4\ldots)$ p. 13-14

$\overline{\underline{\mathrm{I}}}$ ``Caractérisation de $W_\infty$'' p. 1

$\overline{\underline{\mathrm{II}}}$ Perplexités sur les axiomes et les notions (de $W$-fibration notamment) p. 18-22

$\overline{\underline{\mathrm{III}}}$ Morphismes $W$-propres, $W$-lisses, $W$-fibrants p. 22-52

$\overline{\underline{\mathrm{IV}}}$ Digression sur l'axiome $W(4)$ p. 53-58

$\overline{\underline{\mathrm{V}}}$ Sommes amalgamées et carrés \inserted{\unsure{homotopiquement}} \unsure{cocartésiens} dans $(\mathrm{Cat})$ \inserted{(cf notamment th. 4 (p. 145) et th. \unsure{6|5} (p. \unsure{157|177}))} p. 59-169

$\overline{\underline{\mathrm{VI}}}$ \inserted{\S\S\,1,2} \deleted{Les \unsure{digressions|blessures} \unsure{ici}} Critère de $W_\infty$-asphéricité \deleted{\unsure{en}} en termes de la $0$-connexité des $\Omega^{\unsure{i}}(X,\unsure{x|n})$. \unsure{Suite}. p. 171-\deleted{\unsure{191}}190

\deleted{Digression sur la philosophie des \unsure{fibres} homotopiques etc \textemdash\ p. 191-26\unsure{4|9}}

\deleted{Cofibres $W$-homotopiques dans $\mathrm{Cat}$ : p. 227-239}

\deleted{Retour sur les axiomes $W(7)$, $W(7\mathrm{bis})$ \enspace 265-275}

\unsure{Suite|Sortie} des carrés $W$-cartésiens et cocartésiens \enspace 258-264

\S\,3 Sur la philosophie des fibres homotopiques (\unsure{divagations} initiales) \enspace 191-207

\S\,4 Fibres $W$-homotopiques (1) (notamment a) b) c) p. 215, 218) \enspace 207-218

\S\,5 $W$-types d'homotopie relatifs sur localiseur \unsure{constant} \enspace 218-227

\S\,6 Cofibres $W$-homotopiques dans $\mathrm{Cat}$ \enspace 227-240

\S\,7 Fibres $W$-homotopiques (2) ($+$lemme du \unsure{foncteur|fondeur} adjoint) (Théorie des carrés $W$-cartésiens, p. 258-264) \enspace 240-264

\S\,8 Retour \unsure{sur} les axiomes $W(7)$, $W(7\mathrm{bis})$ \enspace 265-275

\uline{Appendices :} $\overline{\underline{\mathrm{I}}}$) Contre-exemples \deleted{\unsure{réfutés|réfutatifs}} : la notion de flèches tot. $W$-asph : \hfill A1-A16

\diagram{\begin{tikzcd}[ampersand replacement=\&,column sep=small,row sep=tiny]
\& \& \bullet\\
\Psi={} \& \bullet \arrow[ur,dash] \arrow[dr,dash] \& \& \text{n'est \uline{pas} \unsure{tot.|LA} $W_2$-asph.}\\
\& \& \bullet
\end{tikzcd}}

$\overline{\underline{\mathrm{II}}}$) Recollement d'un ouvert et d'un fermé comme \inserted{\unsure{$W_\infty$-|théo-}} \unsure{somme} \unsure{amalgamée} (conjectural) \hfill A17-A18

% archive sheet number: 51
\scanpage{52}
\mspage{1}

I) Caractérisation de $W_\infty$ \hfill 19.11.90

Je voudrais caractériser $W_\infty$ en tant que partie de $\mathrm{Fl}(\mathrm{Cat})$, par ses propriétés comme localiseur, et sans faire appel à la connaissance de notions \unsure{utilisant|utilisées} \unsure{déjà|déf.} $W_\infty$, comme le proposait \unsure{Artin|etc.}. Je voudrais tout au moins des propriétés sur un localiseur $W$, familières pour $W_\infty$, qui \inserted{soient si \unsure{fortes} qu'elles} impliquent $W\subset W_\infty$. Bien sûr, dans les démonstrations je me permet d'utiliser tout ce que je sais sur $W_\infty$.

Bien sûr, il faudra au moins une forme de saturation, peut-être très faible, je présume que \unsure{suivant} suffit (non, j'ai dû \inserted{\unsure{le renforcer en (1)}})

% slanted
\margin{$W(1)$ Axiome de saturation \unsure{Ens}}

\textcircled{1} a) Les \deleted{\unsure{isomorphismes}}\inserted{identités} sont dans $W$ \deleted{\unsure{et de}} b) si $f,g$ composables et deux parmi $f,g,gf$ sont dans $W$, la troisième aussi, c) si $fg\in W$, $gf\in W$, alors $f,g\in W$.

% slanted
\margin{\uline{NB}}

Une \deleted{\unsure{axiome}} propriété particulièrement cruciale, qui déjà distingue $W_\infty$ de l'ens des homotopismes (notion trop stricte à beaucoup d'égards, est l'axiome que j'ai appelé L4 bis :

% slanted
\margin{$W(2)$ Axiome de localisation}

\textcircled{2} Soit $f:X\unsure{\text{---}|\longrightarrow}Y$ une flèche dans $\mathrm{Cat}_{/S}$, telle que $\forall s\in\mathrm{Ob}\,S$ $f_s:X_{/s}\longrightarrow Y_{/s}$ soit dans $W$ (on dira que $f$ est \deleted{\unsure{loc}} dans $W$ ``localement sur $S$''. Alors $f\in W$.

\deleted{\uline{Proposition}} \deleted{\unsure{Démon}} Le cas $Y=S$, cela donne

\uline{Corollaire} : Si $f:X\longrightarrow Y$ est $W$-asphérique i.e. $y\backslash X\longrightarrow y\backslash Y$ \deleted{\unsure{dans}} dans $W$ pour tout $y\in\mathrm{Ob}\,Y$, alors $f\in W$.

% slanted
\margin{On peut considérer \deleted{\unsure{cette}} \unsure{propriété} du cor. comme une version affaiblie $W(2\,\mathrm{bis})$ de $W(2)$ ci-dessus. Dans la \unsure{prop. 1} on se \unsure{borne} à $W(2\,\mathrm{bis})$, \unsure{au lieu de} $W(2)$ \deleted{\unsure{il y a dans ce cas a) b) c) sont équivalentes}}}

\uline{Proposition 1} Si $W$ satisfait \textcircled{1} et \textcircled{2}, \deleted{\unsure{deux}} \deleted{\unsure{assertions}} les propriétés suivantes \deleted{\unsure{sont équivalentes}} sont équivalentes :

a) Pour tout $X$ dans $\mathrm{Cat}$, $X\times\Delta^1\mathrel{\mathop{\longrightarrow}\limits^{\mathrm{pr}_1^X}}X$ est dans $W$ (i.e. $\Delta^1\longrightarrow e$ ``universellement dans $W$'').

b) Si $X$ dans $\mathrm{Cat}$ a un objet final, alors $X$ est $W$-asph. i.e. $X\longrightarrow e$ est dans $W$.

c) Tout homotopisme est dans $W$. \deleted{Alors $b\Rightarrow a,\ c\Rightarrow a$, et si les \unsure{flèches} sont \unsure{prises} \unsure{ensemble}, $a\Rightarrow b$, $a\Rightarrow c$, \unsure{b}.}

\uline{Dém} b) $\Rightarrow$ a), car pour voir que $\mathrm{pr}_1^X\in W$, en \unsure{appliquant} le corollaire ci-dessus, il faut prouver que $\forall x\in X$

\[
(X\times\Delta^1)_{/x}\longrightarrow X_{/x}\tag*{(\unsure{*|x})}
\]

est dans $W$. Or
\[
(X\times\Delta^1)_{/x}\simeq X_{/x}\times\Delta^1,
\]

\diagram{\begin{tikzcd}[ampersand replacement=\&,column sep=small,row sep=small]
\Delta^1 \arrow[d,dash] \& \deleted{\unsure{X}}X\times\Delta^1 \arrow[l] \arrow[d] \& X_{/x}\times\Delta^1 \arrow[l] \\
e \& X \arrow[l] \& X_{/x} \arrow[l,dash]
\end{tikzcd}}

et les deux catégories $X_{/x}$, $X_{/x}\times\Delta^1$ sont $W$-asph. car elles ont un objt final, donc on est \unsure{dans} $W$ par saturation \textcircled{1} : \deleted{car $X\times\Delta^1\longrightarrow X$ est \unsure{un homotopisme}.}

% slanted
\margin{on utilise \textcircled{1} (en sa forme faible, sans la condition de saturation c)}

% slanted
\margin{NB Seul $b\Rightarrow a$ utilise $W(2\,\mathrm{bis})$, les deux implications $a\Rightarrow c\Rightarrow b$ sont valables sous la seule hyp\textsuperscript{se} \textcircled{1} ; \deleted{$\unsure{c}\Rightarrow\unsure{a}$} $a\Rightarrow c$.}

Je rappelle que
\[
f:X\longrightarrow Y
\]

% archive sheet number: 52
\scanpage{53}
\mspage{2}

est un homotopisme s'il existe \deleted{\unsure{$g$}} 
\[
g:Y\longrightarrow X
\]

telle que
\[
gf\sim\mathrm{id}_X,\quad fg\sim\mathrm{id}_Y \tag*{(\unsure{*|H})}
\]
où la relation $\sim$ désigne l'homotopie dans $\underline{\mathrm{Hom}}(X,X)$ et $\underline{\mathrm{Hom}}(Y,Y)$ respectivement (définie par chemins allant d'un élément à l'autre). Prouvons le

\uline{Lemme} Supposons que $W$ satisfait \textcircled{1} et le a) de la prop. i.e $X\times\Delta^1\longrightarrow X$ dans $W$ pour tt $X$ dans $\mathrm{Cat}$. Alors si \deleted{\unsure{$f,g$}} $X,Y$ dans $\mathrm{Cat}$ et $f,g$ dans $\underline{\mathrm{Hom}}(X,Y)$ sont homotopes, on a
\[
\mathrm{hot}_W(f)=\mathrm{hot}_W(g).
\]
De plus, \deleted{\unsure{on}} $f\in W$ ssi $g\in W$.

Il suffit de le prouver quand il existe un \unsure{hom|lien} $f\longrightarrow g$. Cela peut s'interpréter par l'existence d'une flèche
\[
X\times\Delta^1\xrightarrow{F}Y
\]

dont les restrictions à $X\times\{0\}$ et $X\times\{1\}$ respectivement \deleted{sont}\inserted{soient} $f,g$. \unsure{Pour} prouver que $\mathrm{hot}_W(f)=\mathrm{hot}_W(g)$, il suffit de prouver que
\[
i_0,i_1:X\unsure{\text{---}|\longrightarrow}X\times\Delta^1
\qquad x\longmapsto(x,0),\ x\longmapsto(x,1)
\]
sont égales dans $\mathrm{Hot}_W$, puisque
\[
f=F\circ i_0,\quad g=F\circ i_1.
\]

Or cela résulte du fait que ce sont des sections de $X\times\Delta^1$ sur $X$, et que dans $\mathrm{Hot}_W$ la projection $X\times\Delta^1\xrightarrow{p}X$ est iso. (donc $i_0=i_1=p^{-1}$ dans $\mathrm{Hot}_W$).

\inserted{De plus, le diagramme}
\diagram{\begin{tikzcd}[ampersand replacement=\&,row sep=tiny,column sep=small]
X \arrow[r,"i_0"] \arrow[dr,"f"'] \& X\times\Delta^1 \arrow[d,"F"]\\
\& Y
\end{tikzcd}}
\inserted{où $i_0\in W$, montre que $f\in W$ ssi $F\in W$.}

% slanted
\margin{et de \unsure{même} $g\in W$ ssi $F\in W$, d'où $f\in W\Longleftrightarrow g\in W$.}

Revenons à (*), il vient que dans $\mathrm{Hot}_W$ $gf=\mathrm{id}_X$, $fg=\mathrm{id}_Y$, donc $f$ et $g$ sont des $\mathrm{Hot}_W$-isom. inverses l'un de l'autre. \deleted{Pour \unsure{aller} plus \unsure{loin}} D'autre part, $gf\in W$, $fg\in W$, \deleted{conclure $f\in W$ ceci \unsure{vient} \unsure{du fait} que $W$ est \unsure{fortement} saturé !} puisque $\mathrm{id}_X$, $\mathrm{id}_Y\in W$. Donc on \uline{gagne}, si on a formulé \textcircled{1} de façon assez forte. On a prouvé donc $a\Rightarrow c$.

c) $\Rightarrow$ b) est trivial, puisque si $X$ a un objet final, $X\longrightarrow e$ est un homotopisme.

% archive sheet number: 53
\scanpage{54}
\mspage{3}

On doit bien sûr pour la suite supposer que $W$ satisfait aux conditions équivalentes de la proposition. On prendra p. ex la forme

% slanted
\margin{$W(3)$ Axiome de l'élément final}

\textcircled{3} Si $X$ a un objet final, alors $X$ est $W$-asphérique.

\uline{Corollaire 1} Pour que $f:X\unsure{\text{---}|\longrightarrow}Y$ soit $W$-asphérique, il f. et s. que les $X_{/y}$ soient $W$-asphériques $(y\in\mathrm{Ob}\,Y)$.

Car pour que
\diagram{\begin{tikzcd}[ampersand replacement=\&,column sep=small,row sep=tiny]
X_{/y} \arrow[rr] \arrow[dr,dash] \& \& Y_{/y} \arrow[dl,dash]\\
\& e
\end{tikzcd}}
soit dans $W$, comme $Y_{/y}$ est $W$-asph \inserted{(objet final)}, il f. et s. que $X_{/y}$ le soit aussi.

\deleted{\uline{\unsure{Proposition}}} \uline{Corollaire 2} Si $X$ dans $\mathrm{Cat}$ a un objet initial, il est $W$-asphérique. (En effet, $X\longrightarrow e$ est un homotopisme, et on applique Prop 1)

% slanted
\margin{C'est donc l'énoncé dual de l'axiome \textcircled{3} de l'objet final.}

Je mettrai pour mémoire

\uline{Corollaire 3} Tout homotopisme dans $\mathrm{Cat}$ est dans $W$. Si $f,g:X\rightrightarrows Y$ \deleted{$\mathrm{Cat}$} dans $\mathrm{Cat}$ sont homotopes, alors elles sont égales dans $\mathrm{Hot}_W$, et $f$ est dans $W$ ssi $g$ l'est.

\uline{Proposition 2} : Soit $f\unsure{\in|:}\deleted{\unsure{W}}X\longrightarrow Y$ dans $\mathrm{Cat}$. Si $f\in W$, $f^\circ\in W^\circ$.

On utilise la dém. de I p. 75, à l'aide de la catégorie \unsure{nommée} $\unsure{\Phi|\mathscr{F}}(X)$ $(\mathrm{Ob}\,\Phi(X)=\mathrm{Fl}(X)$, mais avec une notion de \deleted{\unsure{flèches}}\unsure{flèches} dans $\Phi(X)$ un peu \unsure{particulière|particularisée}). On trouve que $f\in W\Longleftrightarrow\Phi(f)\in W$, et de même $f^\circ\in W\Longleftrightarrow\deleted{\unsure{f^\circ\in W}}\Phi(f)\in W$.

NB. Dans la dém. de la prop 2, comme dans celle de la prop. 1 (pour l'implication b) $\Rightarrow$ a)), on n'a utilisé la propriété de localisation \textcircled{2} que sous la forme faible du corollaire, i.e. \deleted{\unsure{comme}} sous la forme $f$ $W$-asphérique $\Rightarrow f\in W$.

% slanted
\margin{donc si on remplace dans cet énoncé $W(2)$ par $W(\unsure{2}\,\mathrm{bis})$.}

\uline{Corollaire} (de la prop. 2) : Soit $W\subset\mathrm{Fl}(\mathrm{Cat})$ satisfaisant l'axiome \textcircled{1} (saturation). Alors la conjonction des axiomes \textcircled{2} \textcircled{3} équivaut à celle des axiomes duaux \textcircled{2$'$} \textcircled{3$'$}.

Il suffit de prouver que $1,2,3\Longrightarrow\unsure{2'|1'},3'$,

% archive sheet number: 54
\scanpage{55}
\mspage{4}

Pour $2'$, c'est le cor. 2 ci-dessus, \deleted{\unsure{Pour $3'$ c'est}}\inserted{pour $3'$, cela} résulte aussitôt de la prop. 2. (NB Cet énoncé resterait valable, si on avait pris \textcircled{2} sous la forme faible seulement, cf le NB page précédente)

Je dois encore mettre un axiome de la forme L5, L6 (cf $\overline{\underline{\mathrm{VI}}}$, p. 43-44). Mais l'énoncé de ces axiomes utilise la notion de $\mathrm{Hot}$-\unsure{fibrés}, ou de \deleted{\unsure{morphismes}} foncteurs parfaits, i.e. fait appel à la connaissance de $W_\infty$, ce que j'aimerais éviter (si je peux). Je dois donc introduire la notion de $W$-fibration, sur le modèle de ce que j'ai fait dans le contexte des dérivateurs, dans I (cf p. 111-\unsure{126|125}). On introduit la notion de $\mathbb{D}$-fibration à la p\unsure{age} 117 (\unsure{prop} 5), \unsure{quatre|par} conditions équivalentes, dont certaines ont un sens en termes de la donnée de $W$ sans plus. Soit donc \deleted{$f$}

\[
f:X\unsure{\text{---}|\longrightarrow}Y
\]
\deleted{\unsure{a)}} dans $\mathrm{Cat}$, et considérons les conditions

a) Si $Y''\longrightarrow Y'$ dans $\mathrm{Cat}_{/Y}$ est dans $W$, alors $X''\longrightarrow X'$ aussi (le chgt de base par $X\longrightarrow Y$ transforme $W$ dans $W$)

a$'$) Comme a), mais avec $Y'',Y'$ $W$-asphériques.

$a_{\deleted{\unsure{1}}1}$) Comme a), dans le cas particulier où $Y'$ a un objet final \deleted{\unsure{et}} et $Y''$ \deleted{\unsure{celle}} est la \deleted{\unsure{catégorie de}}\inserted{catégorie ponctuelle} \deleted{\unsure{admise}}

$\deleted{\unsure{a}}a_2$) Dual de $a_1$ avec objet initial de $Y'$.

On a \deleted{\unsure{donc}}
\diagram{\begin{tikzcd}[ampersand replacement=\&,row sep=small,column sep=tiny]
\& a\quad W\text{-fibration }\deleted{\text{\unsure{faible}}}\text{\unsure{ultra}forte} \arrow[d,Rightarrow] \&\\
\& a'\quad W\text{-fibration forte} \arrow[dl,Rightarrow] \arrow[dr,Rightarrow] \&\\
W\text{-fibration gauche}\ a_1 \& \& a_2\quad W\text{-fibration droite}
\end{tikzcd}}

et dans le cas $W=W_\infty$ toutes ces conditions sont équivalentes. J'ai \deleted{\unsure{envie}}\inserted{envie} prendre la plus forte de ces conditions, mais je dois me donner des moyens pour trouver des $W$-fibrations en un sens assez fort. Par exemple, j'aurais \unsure{envie} de savoir que si $X$

% archive sheet number: 55
\scanpage{56}
\mspage{5}

est dans $\mathrm{Cat}$, alors $\underline{\mathrm{Ch}}_\infty(X)\unsure{\text{---}|\longrightarrow}X\times X$ ou $\underline{\mathrm{Ch}}(X)\longrightarrow X\times X$ est une $W$-fibration. Dans les deux cas, \deleted{\unsure{je le}} cela (disons) pour $W_\infty$, je le prouve en prouvant que $p$ est à la fois lisse, et propre. Il faudrait donc revoir les notions de lissité et de propreté dans le contexte de la donnée d'un $W$.

Si je me rappelle bien, la propriété a$'$) était quelque chose d'assez aisément vérifiable, il suffisait de supposer que les
\[
X_y\longrightarrow X_{/y}\quad\text{et }\deleted{\unsure{le}}X_{/y}\mathrel{\text{---}}X_{/y'}
\]
(pour $y\in\mathrm{Ob}\,Y$ resp. $u:y\longrightarrow y'$ dans $\mathrm{Fl}\,Y$) sont dans $W$, et que ça reste vrai par tout chgt de base -- ce qui revient \deleted{\unsure{pratiquement}} à la \deleted{\unsure{axiome}} condition $a_1$ sur $f$.

On pourrait donc formuler une \deleted{\unsure{axi}} condition sur $W$ ainsi :

% slanted
\margin{$\unsure{W(4)|W'(4)}$ Axiome des $W$-fibrations (Voir variante suffisante, dans le NB p. 10) Voir \unsure{cascade|série} de variantes \deleted{\unsure{pages}} 13,14 ci-dessous}

\textcircled{4} \inserted{(faible)} Soit $f:X\longrightarrow Y$ dans $\mathrm{Cat}$ satisfaisant la condition de $W$-fibration $a_1)$, ci-dessus (p. 4), alors elle satisfait même : a) i.e le changement de base par $f$ transforme flèches dans $W$ en flèches dans $W$. \inserted{i.e. une $W$-fibration à gauche est une fibr. ultraforte}

Grâce à la prop. 2, il s'ensuit que la condition duale $a_2)$ sur $f$ implique la condition de $W$-fibration forte, donc l'axiome 4, en apparence asymétrique, \deleted{\unsure{équi}} signifie \unsure{aussi|ainsi} que les \inserted{quatre} conditions de fibration sur $f$ sont équivalentes, il est donc en fait symétrique. \deleted{\unsure{Mai}}

\uline{NB} Je ne vois pas, sans axiomes supplémentaires sur $W$, comment passer de $a_1$ à $a'$ seulement. La difficulté, c'est de montrer, moyennant $a_1$ \unsure{seulement}, que \deleted{\unsure{que}} si $Y$ est $W$-asphérique, l'inclusion $X_y\longrightarrow X$ (où $y\in\mathrm{Ob}\,Y$) est dans $W$. Dans le contexte des dérivateurs $\mathbb{D}$, cela se voyait par voie $\mathbb{D}$-cohomologique.

Cet axiome très fort suffit à impliquer que si $X,Y$ sont $W$-fibrés sur $S$, et si $f:X\longrightarrow Y$ est un $S$-morphisme, qui est dans $W$, alors $f_s:X_s\longrightarrow Y_s$ est dans $W$ $\forall s\in S$ ? Montrons déjà au moins l'inverse.

% archive sheet number: 56
\scanpage{57}
\mspage{6}

% slanted
\margin{\uline{sans} axiome \textcircled{4} C'est la première fois qu'on a besoin de $W(2)$ sous sa forme forte, à l'exclusion de $W(2\,\mathrm{bis})$.}

\uline{Proposition 3} Soient $X,Y$ dans $\mathrm{Cat}_{/S}$, $f:X\unsure{\text{---}|\longrightarrow}Y$ une $S$-flèche. Supposons \inserted{$X,Y$ $W$-fibrés} que $\forall s\in S$, $f_s:X_s\longrightarrow Y_s$ soit \deleted{dans} dans $W$, alors $f\in W$.

Il s'impose d'utiliser l'axiome de localisation, donc regarder les $f_{/s}$ induits par $f$
\diagram{\begin{tikzcd}[ampersand replacement=\&,row sep=small,column sep=large]
X_{/s}\arrow[r,"f_{/s}"'] \& Y_{/s}\\
X_s\arrow[u,"i_s^X"]\arrow[r,"f_s"] \& Y_s\arrow[u,"i_s^Y"']
\end{tikzcd}}

\deleted{\unsure{dans les}} condition $a_1$ suffit à impliquer que \deleted{\unsure{les}} $i_s^X$ et $i_s^Y$ sont dans $W$, donc $f_{/s}$ est dans $W$ car $f_s$ l'est. (On n'a pas eu à utiliser l'axiome \textcircled{4}, et il suffisait de supposer que \deleted{\unsure{$f$}}$X,Y$ soient fibrés ``$W$-cohomologiques'' \deleted{\unsure{sur $S$}} (ou \inserted{$W$-fibrés} à g.) sur $S$, ou d'ailleurs, des $W$-fibrés à dr.

De plus :

\uline{Corollaire} \deleted{\unsure{$S$}} Soit $S_0\subset\mathrm{Ob}\,S$ une partie de $S$ telle que $S_0\longrightarrow\pi_0(S)$ surjectif. Alors dans la prop 3 il suffit de supposer $f_s\in W_{\unsure{4|s}}$ pour $s\in S_0$.

% slanted
\margin{sans axiome \textcircled{4}}

Il suffit de voir que si $u:s\longrightarrow s'$ est une flèche dans $\unsure{s|S}$, alors $f_s\in W\Longleftrightarrow f_{s'}\in W$. On utilise le fait que $f_s\in W\Longleftrightarrow f_{/s}\in W$, et de même pour $f_{s'}$. Donc à voir que $f_{/s}\in W\Longleftrightarrow f_{/s'}\in W$. Or on a

\diagram{\begin{tikzcd}[ampersand replacement=\&,row sep=large,column sep=large]
X_{/s}\arrow[r,"f_{/s}"]\arrow[d,"X_{/u}"'] \& Y_{/s}\arrow[d,"X_{/u}"]\\
X_{/s'}\arrow[r,"f_{/s'}"] \& Y_{/s'}
\end{tikzcd}}
où $X_{/u}$ et $Y_{/u}$ sont dans $W$, d'où le résultat.

Mais revenons à la réciproque : si $f\in W$, alors a-t-on $f_s\in W_s$ ? Il y a un argument dans ce sens dans $\overline{\underline{\mathrm{VI}}}$ p. 23, que je vais reprendre. L'idée \deleted{c'est de} est celle-ci. On considère la comp. conn. $Y_0$ de $Y$ contenant $s$, et on \deleted{\unsure{est}} construit un

% archive sheet number: 57
\scanpage{58}
\mspage{7}

$W$-fibré (au \uline{sens fort} a))
\[
S'\longrightarrow S_0\subset S
\]

avec $S'$ $W$-\uline{asphérique}. Comme $X\unsure{\text{---}|\longrightarrow}Y$ sur $S$ est dans $W$, il en sera donc de même de $X'\mathrel{\text{---}}Y'$ sur $S'$. \deleted{\unsure{Cela}} Comme les fibres de $X',Y'$ en $s'$ s'identifient à celles de $X,Y$ en l'image de $s'$ dans $S$, on est ramené au cas \inserted{où} $S$ \deleted{\unsure{asp}} est $W$-asphérique. Mais alors on a

\diagram{\begin{tikzcd}[ampersand replacement=\&,column sep=large]
X\arrow[r,"f"] \& Y\\
X_s\arrow[u,"i_s^X"]\arrow[r,"f_s"'] \& Y_s\arrow[u,"i_s^Y"']
\end{tikzcd}}

, où $i_s^X,i_s^Y$ sont dans $W$, car déduites de $\{s\}\longrightarrow S$ (qui est dans $W$, $S$ étant $W$-asphérique) par les chgt. de base $W$-fibrants $X\longrightarrow S$, $Y\longrightarrow S$. Donc $f\in W\Longrightarrow f_s\in W$.

NB On n'a utilisé la propriété de fibration pour $X,Y$ \unsure{au} sens a$'$).\textsuperscript{$*$}

% slanted
\margin{Et même, dans a$'$) on pourrait \deleted{\unsure{exiger}} que $Y''$ \unsure{soit} $=e$, $Y'$ $W_\omega$-asphérique, \uline{\unsure{page} 14}, $W_\omega\subset W\cap W_\infty$ ?}

Il faut donc \deleted{\unsure{voir}} prouver, si on peut, le

\uline{Lemme} (?) Soit $S$ un objet $0$-connexe de $\mathrm{Cat}$. Alors il existe un objet $W$-asphérique $S'$ de $\mathrm{Cat}$, et une \deleted{\unsure{$\mathrm{Hot}$}}-fibration \inserted{(au sens \uline{fort} de a) page 4)} $S'\longrightarrow S$.

% slanted
\margin{Est prouvé moyennant l'axiome \textcircled{4} ($+$ des petites vérifications \unsure{ultérieures})}

L'idée naturelle, c'est de prendre, pour $s_0\unsure{\subset|\in}\mathrm{Ob}\,S$ fixé,
\[
S'=\underline{\mathrm{Ch}}(S;s_0,-)\quad\text{ou}\quad S'=\underline{\mathrm{Ch}}_\infty(S;s_0,-),
\]
et la flèche
\[
\beta:c\longmapsto\mathrm{but}\,c:S'\longrightarrow S.
\]

On sait que $S'$ est $W_\infty$-asphérique et que $\beta$ est $W_\infty$-fibrante (et même parfaite). Il faut voir que $S'$ est aussi $W$ asphérique, et $\beta$ $W$-fibrante. Il faut donc revoir la démonstration dans le cas de $W_\infty$, et prouver qu'elle se \unsure{transfère|transforme} à $W$,

% archive sheet number: 58
\scanpage{59}
\mspage{8}
% archive sheet number: 59

moyennant les axiomes qu'il faut sur $W$. J'ai envie de regarder d'abord le choix $S'=\underline{\mathrm{Ch}}(S;s_0,-)$, sans passage à une limite inductive qui m'obligerait à mettre un axiome de limites sur $W$. Les \deleted{\unsure{théorie}} propriétés homotopiques délicates et essentielles de $\underline{\mathrm{Ch}}(S)$ sont établies dans $\overline{\underline{\mathrm{VII}}}$, p. 157-169, savoir que dans le diagramme

\diagram{\begin{tikzcd}[ampersand replacement=\&]
S\arrow[r,"i_S",dash] \& \underline{\mathrm{Ch}}(S)\arrow[r,"p_S"] \& S\times S
\end{tikzcd}}
, $i_S\in W_\infty$ et $p_S$ est une \deleted{\unsure{$W$}}$\unsure{\mathrm{Hot}|W_\infty}$-fibration (en fait, \uline{propre} et \uline{lisse}\inserted{ pour $W_\infty$}). Le fait que $i_S\in W$ implique que les \unsure{foncteurs} source et but

% slanted
\margin{Dans l'énoncé des ``propriétés ess.'' de $\underline{\mathrm{Ch}}(S)$, on peut travailler avec $W_\omega$, $i_S\in$ la plus petite partie de $\mathrm{Fl}(\mathrm{Cat})$ satisfaisant les axiomes $W(1)$, $W(2\,\mathrm{bis})$, $W(3)$ (cf \unsure{p.ge|page} 14), donc $W_\omega\subset W\cap W_\infty$}

\diagram{\begin{tikzcd}[ampersand replacement=\&,column sep=large]
\underline{\mathrm{Ch}}(S)\arrow[r,shift left=.6ex,"\sigma",dash]\arrow[r,shift right=.6ex,"\beta"',dash] \& S
\end{tikzcd}}
\[
\sigma=\mathrm{pr}_1\circ p_S,\quad \beta=\mathrm{pr}_2\circ p_S
\]

sont dans $W_\infty$ (puisque $\sigma i_S=\beta i_S=\mathrm{id}_S$). Comme \deleted{c'est une} sont des $\mathrm{Hot}$-fibrations, il en résulte que les fibres $\underline{\mathrm{Ch}}(S;s,-)$ ou $\underline{\mathrm{Ch}}(S;-,s)$ sont $\unsure{\mathrm{Hot}}$-asphériques. Mais même en admettant que $\sigma,\beta$ sont dans $W$ et sont des \deleted{\unsure{$\mathrm{Hot}$}}$W$-fibrations, je ne peux en conclure \inserted{tout de suite} que les fibres sont $W$-asphériques -- il y aurait cercle vicieux ! Il me faut donc une démonstration directe.

Considérons le foncteur
\diagram{\begin{tikzcd}[ampersand replacement=\&,row sep=small]
\underline{\mathrm{Ch}}(S;s,-)\arrow[d,"p"']\\
\underline{\mathrm{Ty}}
\end{tikzcd}}

, c'est un foncteur $\mathrm{Cat}$-fibrant, donc $W$-lisse (je \deleted{devrai} \unsure{développer} la notion de $W$-lissité \unsure{bien} sûr, \deleted{\unsure{je dois}} \unsure{j'utiliserai|j'admettrai} pour le moment les propriétés formelles familières, et les justifierai par la suite). A cause de la $W$-lissité, pour prouver $p\in W$, il suffit de

\margin{\deleted{\unsure{$f$}}$f:X\unsure{\text{---}|\longrightarrow}Y$ est $W$-lisse si les $X_y\longrightarrow{}_{y\backslash}X$ sont $W$-asphériques}
\scanpage{60}
\mspage{9}
% archive sheet number: 60

% slanted
\margin{on utilise ici $W(2\,\mathrm{bis})$, pas $W(2)$}

prouver que les fibres de $p$ sont $W$-asphériques. Or les fibres sont les $Z_\tau=\underline{\mathrm{Ch}}_\tau(S;s,-)$, qui sont des catégories \uline{contractiles} (i.e. $Z_\tau\unsure{\text{---}|\longrightarrow}e$ est un homotopisme), donc $W$-asphériques par \textcircled{3} (cf \unsure{cor}. 3 page 3).

D'autre part, \deleted{voyons que} si on admet que $\underline{\mathrm{Ch}}(S)\xrightarrow{p_S}S\times S$ est $W$-lisse (et $W$-propre)\deleted{\unsure{et}}, donc $\underline{\mathrm{Ch}}(S)\xrightarrow{\sigma_S}S$ aussi, étant à fibres $W$-asphériques comme on vient de voir, il s'ensuit que $\sigma_S\in W$, d'où $S\xrightarrow{i_S}\underline{\mathrm{Ch}}(S)$ aussi dans $W$ (puisque $\deleted{\unsure{p}}\sigma_S\,i_S=\mathrm{id}_S$).

% slanted
\margin{\unsure{id}}

Donc il reste à prouver seulement que
\[
\underline{\mathrm{Ch}}(S)\xrightarrow{p_S}S\times S \tag*{(a)}
\]
est $W$-\deleted{\unsure{lisse}}propre et $W$-lisse. Il en résultera que
\[
\underline{\mathrm{Ch}}(S;s_0,-)\xrightarrow{\beta}S
\]
est lui aussi $W$-propre et $W$-lisse, car déduit de $p_S$ par le chgt de base
\[
S\longrightarrow S\times S\qquad s\longmapsto(s_0,s).
\]

Il faut ensuite voir que \deleted{\unsure{$f$}} $W$-lisse et $W$-propre implique la propriété de $W$-fibration \uline{sous sa forme la plus forte} a) (alors que je prévois qu'a priori on ne l'aura que sous les formes faibles $a_1,a_2$ -- pas même sous la forme $a'$). C'est ici donc que l'axiome \textcircled{4} des $W$-fibrations sera \unsure{appréciable|apprécié} !

% slanted
\margin{\unsure{sous} la forme affaiblie $\unsure{W(4a'')|4a'')}$ (p. 13)}

Voici l'argument pour la lissité de \unsure{(a)|cas}, dans loc. cit p. 161 ff. Il faut prouver que \deleted{la catégorie}\inserted{pour} \deleted{\unsure{$c\in\mathrm{Ob}\,\underline{\mathrm{Ch}}$}} $c\in\mathrm{Ob}\,\underline{\mathrm{Ch}}(S)$, \deleted{\unsure{$(a,b)\in\mathrm{Ob}$}} et une flèche $u\in\mathrm{Fl}(S\times S)$, une certaine catégorie du ``relèvement de $u$ de source $c$'', $C=C(c,u,\unsure{v|u})$, est $W$-asphérique. Pour ceci, on la compare à une catégorie \inserted{$C'$} à objet final, donc \inserted{contractile et à fortiori} $W$-asphérique, à l'aide de \deleted{\unsure{du}} foncteurs

% slanted
\margin{En fait, on prouve que $C(c,u,v)$ est \uline{contractile}, donc elle est $W$-asphérique}

\diagram{\begin{tikzcd}[ampersand replacement=\&,column sep=large]
C\arrow[r,shift left=.65ex,"\varphi",dash] \& C'\arrow[l,shift left=.65ex,"\psi"]
\end{tikzcd}}
, et des homotopies $\varphi\psi\sim\mathrm{id}_{C'}$ (donc, puisque $C'$ est contractile)
\scanpage{61}
\mspage{10.}

et $\psi\varphi\sim\mathrm{id}_C$ (on trouve même une flèche $\psi\varphi\unsure{\text{---}|\longrightarrow}\mathrm{id}_C$), Ainsi $\varphi$ est un homotopisme, donc une $W$-équivalence (cor. 3 p. 3), donc $C$ est $W$-asphérique, $C'$ l'étant.

La propreté de $p_S$ équivaut à la lissité de $p_{(S^\circ)}$, elle est donc prouvée du même coup.

Ainsi, le lemme de la page 7 est prouvé, moyennant l'axiome \textcircled{4} qui assure qu'une flèche qui est à la fois $W$-lisse et $W$-propre, satisfait à la condition très forte \inserted{a)} des $W$-fibrations.

\margin{\uline{NB}}

En fait, il suffirait, au lieu de l'axiome \textcircled{4} sous \deleted{\unsure{sa}} forme forte, de supposer qu'une $f:X\longrightarrow Y$ \deleted{qui} est à la fois ``homotopiquement lisse'' et ``homotopiquement propre'', i.e. les fibres des $X_{/x}\longrightarrow Y_{/y}$ et des $ {}_{x\backslash}X\longrightarrow{}_{y\backslash}Y$ sont \unsure{contractiles|contractibles} (donc $W$-asphériques pour tout $W$ satisfaisant \textcircled{1}\textcircled{2}\textcircled{3}), soit $W$-fibrante forte. Énonçons quand même ce que nous avons prouvé.

% slanted
\margin{pour \textcircled{2}, il \unsure{le} suffit sous la forme $W(2\,\mathrm{bis})$}

\uline{Proposition 4} Si $W$ satisfait \textcircled{1}\textcircled{2}\textcircled{3} (nous \unsure{les supposons|le supposons} tjrs par la suite) plus \textcircled{4} (sous une forme un peu affaiblie suffirait, cf remarque précédente), on a ceci : Soit $f:X\longrightarrow Y$ une flèche dans $\mathrm{Cat}_{/S}$, avec $X,Y$ des $W$-fibrés sur $S$. \deleted{Alors} Si $f\in W$, alors $f_s:X_s\longrightarrow Y_s$ est dans $W$ pour tout $s\in S$.

Je me demande si j'en sais assez pour pouvoir prouver
\[
W\subset W_\infty.
\]

Il faut quand même éliminer \unsure{du|le} cas \unsure{canularesque|cauchemardesque}, \unsure{horreur|hmmm} $W=\mathrm{Fl}(\mathrm{Cat})$, donc mettons l'axiome

% slanted
\margin{$W(5)$ Axiome de connexité. Voir à ce sujet les cor. page 22}

\textcircled{5} Si $f:X\longrightarrow Y$ est dans $W$, $\pi_0(f):\pi_0(X)\longrightarrow\pi_0(Y)$ bijectif.

\deleted{Essayons de montrer que si $f:X\longrightarrow Y$ est dans $W$, il est dans $W_\infty$. Soit $I=\pi_0(X)\simeq\pi_0(Y)$, considérons}

\deleted{\diagram{\begin{tikzcd}[ampersand replacement=\&,row sep=tiny,column sep=small]
X\arrow[rr,"f"]\arrow[dr] \& \& Y\arrow[dl]\\
\& I
\end{tikzcd}}}

% archive sheet number: 61
\scanpage{62}
\mspage{11}

\deleted{Comme $X$ et $Y$ sont homotopiquement lisses et propres sur $I$, donc des $W$-fibrés sur $I$ par $\textcircled{4}$, on voit que $f_i:X_i\to Y_i$ est dans $W$ pour tt $i\in I$. Donc \unsure{OBS|OPS} $0$-connexes.}

% slanted
\margin{À nouveau, on \unsure{utilise|voit} un \unsure{ingrédient} \unsure{mystérieux}, \unsure{inutile} : \uline{tout \unsure{baigne}}. en fait, \unsure{faut} : il faut $W_\infty$-parfait (i.e. $W_\infty$-lisse et propre) et $W$-parfait [en fait, $W_\infty$-parfait]}

Je factorise donc $f$ en
\[
 X\xrightarrow{i_{\unsure{*|X}}}\bar X\xrightarrow{\bar f}Y,
\]
\deleted{\unsure{cet}} $\bar f$ une $W_\infty$-fibration\inserted{et une $W$-fibration$^*$}, $i_X\in W\cap W_\infty$ --- cela marche fonctoriellement, par la construction de $\overline{\underline{\mathrm{VII}}}$ p.170-171 (du Cartan, Serre), à partir des propriétés de $Y\to\underline{\mathrm{Ch}}(Y)\to Y\times Y$ sans plus.

\deleted{Donc $f\in W\cap W_\infty$} Comme $\bar f i=f\in W$\inserted{et $i\in W$}, on déduit $\bar f\in W$. \deleted{Donc pour prouver $f\in W_\infty$} Comme $i\in W_\infty$, on en \unsure{conclut|déduit} $f\in W_\infty\Longleftrightarrow\bar f\in W_\infty$.

Donc pour prouver que $f$ est dans $W_\infty$, on est ramené au cas où $f$ est une $W$-fibration et une $W_\infty$-fibration. Mais alors la prop. 4 nous montre que les fibres de $f$ sont $W$-asphériques, et il suffit de voir qu'elles sont $W_\infty$-asphériques (prop. $3^*$ appliquée à $W_\infty$),

% slanted
\margin{on n'a utilisé ici que \unsure{cette} propriété, sur les fibres, comme conséq. de $\textcircled{4}$ (ou utiliser 4 bis ou sa forme la plus faible $4\,\mathrm b''$ cf. pages 13, 14.}

\margin{on n'utilise la prop 3 que \unsure{pour|dans} le cas $Y=S$, pour lequel $W(2~\mathrm{bis})$ suffit}

Ainsi on est ramené au cas où $Y=e$, i.e. à prouver que si $X$ dans $\mathrm{Cat}$ est $W$-asphérique, il est $W_\infty$-asphérique. \deleted{L'axi} \deleted{L'axiome de connexité sur $W$ nous assure déjà que $X$ est $0$-connexe.}

Soit $x_0\in\operatorname{Ob}X$, considérons comme \deleted{\unsure{tantôt}}
\[
 E=\underline{\mathrm{Ch}}(X,x_0,-)\xrightarrow{p}X,
\]

on sait que $E$ est $W\cap W_\infty$-asphérique, et que c'est une fibration pour $W$ et pour $W_\infty$. Comme $X$ est $W$-asphérique\inserted{(comme $E$)}, $p\in W$ \deleted{\unsure{est}}, donc \inserted{à nouveau} par la prop. $4^*$ les fibres de $p$ sont $W$-asphériques, en particulier (par $\textcircled{5}$) $0$-connexes. On a donc prouvé le

% slanted
\margin{en fait, la condition $\unsure{W(4\,\mathrm b'')|W4\,\mathrm b''}$ sur $W$}

% archive sheet number: 62
\scanpage{63}
\mspage{12}

\uline{Lemme} Soit $(X,x_0)$ dans $\mathrm{Cat}_{\bullet}$. Si $X$ est $W$-asphérique, alors $\Omega(X,x_0)\overset{\text{déf}}{=}\underline{\mathrm{Ch}}(X;x_0,x_0)$ \deleted{est} l'est aussi (moyennant les axiomes $\textcircled{1}\textcircled{2}\textcircled{3}\textcircled{4}$ sur $W$). Donc (par itération) les $\Omega^n(X,x_0)$ sont $W$-asphériques $(\forall n\in\mathbb N)$.

% slanted
\margin{Utilise $W(1),W(2~\mathrm{bis})$ $W(3),W(4\,\mathrm b'')$ sur $W$}

\uline{Corollaire} Si de plus $W$ satisfait l'axiome $\textcircled{5}$ de connexité, alors $X$ est \deleted{\unsure{donc}} $W_\infty$-asphérique.

En effet, ``on sait'' que cela signifie que tous les $\Omega^n(X,x_0)$ sont $0$-connexes. Or l'axiome de connexité $\textcircled{5}$ assure que $W$-asphérique implique $0$-connexe.

% slanted
\margin{C'est prouvé fort simplement dans le corollaire au ``\unsure{th, 1|th. 1}'', page 190}

Mais le ``on sait'' est un tout petit peu léger$^*$. Pour le justifier, il faudrait, ou bien que je fasse le lien avec la théorie \unsure{$\frac{1}{2}$ simpliciale|ksimpliciale|simpliciale} classique, et prouve que les $\Omega^n$ que je construis sont ``les bons'', i.e. que leur Nerf redonne bien (à isom près dans $\mathrm{Hot}$) les $\Omega^n$ classiques (??). En particulier, que le $\pi_0(\Omega^n(X,x))$ est bien le $\pi_n(X,x_0)$ ``de tout le monde''. Ou bien que je prouve directement que la $0$-connexité des \unsure{\S} $\Omega^n(X,x_0)$ équivaut à la $W_\infty$-asphéricité de $X$, définie comment ? \unsure{grossièrement} : le Artin-Mazur.

C'est donc un travail qu'il me faudra faire \unsure{absolument|obscurément}, Mais ce n'est plus qu'un simple travail de routine, et je \deleted{\unsure{puis}}\inserted{\unsure{peux}} énoncer dès maintenant le

% margin note written sideways
\margin{les hypothèses utiles sur $W$ sont finalement $W(1),W(2~\mathrm{bis}),W(3),W(4\,\mathrm b''),$ et $W\ne\mathrm{Fl}(\mathrm{Cat})$}

% slanted
\margin{Il suffit même de $\textcircled{4}$ sous une forme affaiblie, une des formulations ``faibles'' \unsure{des nos|de mes} \unsure{deux} ``\unsure{absolus|absol.}'' cf le NB.}
% slanted
\margin{\unsure{ci|ce} suffit, cf. \unsure{infra} p. 13, 14.}

\uline{Théorème 1}. Soit $W\subset\mathrm{Fl}(\mathrm{Cat})$ satisfaisant les conditions $\textcircled{1}\textcircled{2}\textcircled{3}\textcircled{4}\textcircled{5}$ (axiomes de saturation, de localisation, de l'objet final, de fibration, de connexité). Alors $W\subset W_\infty$.

\margin{\uline{NB} Au lieu de $W(5)$, il suffit de supposer $W\ne\mathrm{Fl}(\mathrm{Cat})$, cf \deleted{\unsure{th. 2}} p. 20.}

\uline{Corollaire} On peut caractériser $W_\infty$ comme la plus grande partie de $\mathrm{Fl}(\mathrm{Cat})$ satisfaisant aux

% archive sheet number: 63
\scanpage{64}
\mspage{13}

conditions $\textcircled{1}$ à $\textcircled{5}$ précédentes.

% slanted
\margin{cf axiome $W(6)$ page suivante}

Bien sûr, on aimerait un axiome 6 plus satisfaisant, qui joint aux précédents assure que $W=W_\infty$.$^*$ Ou est-ce que ces propriétés assurent déjà que $W=W_\infty$ ? Pour le voir, il faudrait peut-être que j'en sache plus sur la structure des objets $W_\infty$-asphériques ? Si c'était faux, je songe à un axiome de ``descente homotopique'', qui indiquerait qu'un objet $W_\infty$-asphérique est aussi $W$-asphérique.

Je note dès maintenant que si des $W_i$, $i\in I$, satisfont aux axiomes $1^\circ$ à $5^\circ$, il en est de $=$ des $\bigcap W_i$. Donc il y a un plus petit ens. de flèches $W_{\unsure{\omega|\infty}}$ satisfaisant à ces conditions. S'il \unsure{m'est|m'était|était} égal à $W_\infty$, i.e. si les axiomes 1 à 5 ne caractérisent déjà $W_\infty$, ce serait un \deleted{\unsure{résultat}} localiseur vraiment très intéressant !

% slanted
\margin{C'est vrai pour tout ens. d'axiomes dans \unsure{la catégorie|le catalogue} d'axiomes de stabilité}

Je me rends compte que je n'ai utilisé qu'une forme assez faible de l'axiome $\textcircled{4}$ -- essentiellement, la validité de la prop. 4. Je vais résumer différentes variantes de l'axiome de fibration.

$\textcircled{\underline{4}}$ a$_1$) Si $f$\unsure{$:$|$\in$}$X\to Y$ dans $\mathrm{Cat}$ est une $W$-fibration gauche, c'est une $W$-fibration \inserted{ultra} forte

a$_2$) Si $f:X\to Y$ est une $W$-fibration droite, c'est une $W$-fibration \inserted{ultra} forte

\deleted{\unsure{b|d}}a$'$) Si $f:X\to Y$ est \inserted{$W$-parfait, i.e.} à la fois $W$-propre et $W$-lisse$^*$, alors $f$ est une $W$-fibration \inserted{ultra} forte

% slanted
\margin{ce qui implique que c'est une $W$-fibration gauche \uline{et} droite}

a$''$) Si $f:X\to Y$ est à la fois homotopiquement propre et homotopiquement lisse (cf NB page 10),$^{\dagger}$ c'est une $W$-fibration \inserted{ultra} forte. $^{\dagger}$ Il vaut mieux dire : si $f$ est $W_\infty$-parfait, c'est une $W$-fibr. ultra-forte

\deleted{\unsure{d|c}}b) Soit $f:\unsure{\alpha|X}\to Y$ une flèche de $\mathrm{Cat}_{/S}$, avec $X$ et $Y$ des $W$-fibrés sur $S$. Si $f\in W$, alors $\forall s\in S$, $f_s:X_s\unsure{\mathrel{-}|\to}Y_s$ est $\in W$.

\deleted{\unsure{d|c}}b$'$) Soit $f:X\to Y$ une flèche de $\mathrm{Cat}_{/S}$, qui soit \deleted{à la fois} $W$-parfaite et $W_\infty$-parfaite, \deleted{\unsure{et si}} alors $f\in\unsure{V|W}$ implique que les fibres de $f$ sont $W$-asphériques. (Il vaut mieux dire : si $f$ est $W_\infty$-parfait, il est une $W$-fib. ultra-forte)

% archive sheet number: 64
\scanpage{65}
\mspage{14}

\deleted{\unsure{d|b}}\unsure{b|d}$''$) Soit $W_\omega\deleted{\unsure{\in|=}}\subset W\cap W_\infty$ la plus petite partie de $\mathrm{Fl}(\mathrm{Cat})$ satisfaisant les conditions $\textcircled{1}\textcircled{2}\textcircled{3}{}^*$ (axiome de saturation, de localisation, de l'objet final). Soit $f:X\to Y$ une flèche de $\mathrm{Cat}$ qui soit $W_\omega$-parfaite. Si $f\in W$, alors les fibres de $f$ sont $W_\omega$-asphériques.

% slanted
\margin{Je vais même prendre $\text{\textcircled{\scriptsize 2 bis}}$ au lieu de $\textcircled{2}$, pour trouver $W_\omega$ le plus petit possible !}

Entre ces propriétés de $W$, on a les \unsure{relations|implications} tautologiques (sauf $a''\Rightarrow b)$, qui est \unsure{exact.} la prop. 4)

\deleted{\diagram{\begin{tikzcd}[ampersand replacement=\&,row sep=small,column sep=tiny]
\unsure{a_1|a}\arrow[dr,Rightarrow] \& \unsure{a|a_1}\arrow[d,Rightarrow]\arrow[r,Leftrightarrow] \& \unsure{a_2|a}\arrow[dl,Rightarrow]\\
\& b\arrow[d,Rightarrow]\\
\& c\arrow[d,Rightarrow]\\
\& d\arrow[d,Rightarrow]\\
\& d'\arrow[d,Rightarrow]\\
\& d''.
\end{tikzcd}}}

\diagram{\begin{tikzcd}[ampersand replacement=\&,row sep=small]
a=(a_1\Leftrightarrow a_2)\arrow[d,Rightarrow]\\
a'\arrow[d,Rightarrow]\\
a''\arrow[d,Rightarrow]\\
b\arrow[d,Rightarrow]\\
b'\arrow[d,Rightarrow]\\
b''
\end{tikzcd}}

\margin{$\left.\vphantom{\begin{matrix}a\\a'\\a''\end{matrix}}\right\}$ En fait, toutes ces conditions sont équivalentes.}

% slanted
\margin{Je reviens sur les diverses variantes de $W(4)$ plus bas, p. \unsure{53-60|53-66|53-56}}

Eh bien, dans le lemme page 12, il suffisait de supposer que $W$ satisfait à la version la plus faible $d'')$ de l'axiome de fibration (sept versions !). En effet, les fibrations $\bar X\unsure{\mathrel{-}|\to}Y$ et $E\to X$ utilisées dans la démonstration du lemme sont $W_\omega$-parfaites, et il suffisait de savoir que si elles sont $\in W$, alors les fibres sont $W$-asphériques.

Mais je ne pense pas que même la version la plus forte $4\,\mathrm a)$ de \deleted{\unsure{fibra}} l'axiome de fibration, jointe aux trois axiomes précédents\inserted{et à 5 (connexité)}, implique $W=W_\infty$. Il faudrait tout au moins y joindre l'axiome.

% slanted
\margin{$W(6)$ Axiome des limites filtrantes}

$\textcircled{6}$ Soit $X=\underset{i}{\underline{\lim}}X_i$ une limite inductive filtrante dans $\mathrm{Cat}$. Si les $X_i$ sont $W$-asphériques, $X$ aussi.

Ceci devrait être équivalent à la propriété, en apparence plus forte, en termes de \deleted{\unsure{toutes}} flèches $\in W$ :
\uline{Proposition \unsure{5}} Supposons qu'on ait$^{\unsure{1|\wedge}}$ $\textcircled{6}$, (en plus de $\textcircled{1}$ à $\textcircled{4}{}^*$),

% slanted
\margin{$\textcircled{4}$ sous la forme la plus faible \deleted{\unsure{$4\,\mathrm b$}} $\text{\textcircled{\scriptsize $4\,\mathrm d''$}}$}

% archive sheet number: 65
